\begin{thebibliography}{99}
\bibitem[Higgott(2024)]{thesis} Higgott, O. (2024). \emph{Practical and Efficient Quantum Error Correction}. Ph.D. dissertation, University College London. Submitted October 1, 2023. \url{https://discovery.ucl.ac.uk/id/eprint/10186408/}.
\bibitem[Higgott and Breuckmann(2023)]{construction} Higgott, O., and Breuckmann, N. P. (2023). Constructions and performance of hyperbolic and semi-hyperbolic Floquet codes. arXiv:2308.03750v1. \url{https://arxiv.org/abs/2308.03750v1}. Source data: \url{https://github.com/oscarhiggott/hyperbolic-floquet-data}.
\bibitem[Gidney et al.(2022)]{honeycomb} Gidney, C., Newman, M., and McEwen, M. (2022). Benchmarking the Planar Honeycomb Code. arXiv:2202.11845v3. \url{https://arxiv.org/abs/2202.11845v3}. Generator: \url{https://github.com/Strilanc/honeycomb-boundaries}.
\bibitem[Gidney(2021)]{stim} Gidney, C. (2021). Stim: a fast stabilizer circuit simulator. \emph{Quantum}, 5, 497. \url{https://doi.org/10.22331/q-2021-07-06-497}. Software versions 1.16.0 and 1.12.0: \url{https://github.com/quantumlib/Stim}.
\bibitem[Higgott(2021)]{matching} Higgott, O. (2021). PyMatching: A Python package for decoding quantum codes with minimum-weight perfect matching. arXiv:2105.13082v2. \url{https://arxiv.org/abs/2105.13082v2}. Software versions 2.4.0 and 2.1.0: \url{https://github.com/oscarhiggott/PyMatching}.
\bibitem[Holm(1979)]{multiple} Holm, S. (1979). A simple sequentially rejective multiple test procedure. \emph{Scandinavian Journal of Statistics}, 6(2), 65--70. \url{https://www.jstor.org/stable/4615733}. We use the conservative single-step Bonferroni bound, not the step-down refinement; bibliographic metadata and abstract were accessible, not the complete original article.
\bibitem[Software projects(2026)]{software} Software projects. Version-specific packages used for implementation: Python 3.12.14 (\url{https://www.python.org/}); NumPy 2.5.4 (\url{https://github.com/numpy/numpy}); SciPy 1.18.1 (\url{https://github.com/scipy/scipy}); NetworkX 3.7 (\url{https://github.com/networkx/networkx}); Matplotlib 3.11.2 (\url{https://github.com/matplotlib/matplotlib}). These packages support numerical arrays, normal quantiles, independent graph matching, and plotting respectively. Pytest 9.2.0 (\url{https://github.com/pytest-dev/pytest}) runs the checks; pdfTeX 1.40.26 (\url{https://www.tug.org/texlive/}) renders the paper. The historical environment uses Python 3.10.22, NumPy 1.24.4, SciPy 1.10.1, NetworkX 2.8.8, and Matplotlib 3.7.3, with its dependency lock included.
\bibitem[Harris et al.(2020)]{numpy} Harris, C. R., et al. (2020). Array programming with NumPy. \emph{Nature}, 585, 357--362. \url{https://doi.org/10.1038/s41586-020-2649-2}.
\bibitem[Virtanen et al.(2020)]{scipy} Virtanen, P., et al. (2020). SciPy 1.0: Fundamental algorithms for scientific computing in Python. \emph{Nature Methods}, 17, 261--272. \url{https://doi.org/10.1038/s41592-019-0686-2}.
\bibitem[Hagberg et al.(2008)]{networkx} Hagberg, A. A., Schult, D. A., and Swart, P. J. (2008). Exploring network structure, dynamics, and function using NetworkX. \emph{Proceedings of the 7th Python in Science Conference}, 11--15. \url{https://conference.scipy.org/proceedings/SciPy2008/paper_2/}.
\bibitem[Hunter(2007)]{matplotlib} Hunter, J. D. (2007). Matplotlib: A 2D graphics environment. \emph{Computing in Science \& Engineering}, 9(3), 90--95. \url{https://doi.org/10.1109/MCSE.2007.55}.
\end{thebibliography}